% !TEX root = ../../main.tex
% C4.4 — Sơ đồ use case tổng thể
% Nguồn: usecase_detail.md (dòng Actors + các câu "chuyển sang UC-xx"). Hình H1: SV vẽ theo figures/specs/H1-use-case.md.
% Quan hệ «extend»: UC-10→UC-09, UC-11→UC-10, UC-13→UC-12, UC-14→UC-13. Không vẽ «include» tới Đăng nhập (tiền điều kiện).
\section{Sơ đồ use case tổng thể}
\label{sec:4.4}

Các yêu cầu chức năng được tổ chức thành 15 use case, thể hiện cùng tác nhân và quan hệ trong Hình~\ref{fig:usecase} và liệt kê kèm mục đích trong Bảng~\ref{tab:usecases}.

\begin{figure}[H]
\centering
\includegraphics[width=\textwidth]{H1-use-case.png}
\caption{Sơ đồ use case tổng thể của hệ thống}
\label{fig:usecase}
\end{figure}

Sơ đồ theo bốn lựa chọn sau:
\begin{itemize}
    \item \textbf{Tác nhân trừu tượng \emph{Người dùng}.} Đăng nhập (UC-01) và đăng xuất (UC-15) được nối với tác nhân trừu tượng \emph{Người dùng}, mà ba tác nhân cụ thể tổng quát hóa về.
    \item \textbf{Phân chia theo vai trò.} Quản trị viên có UC-02 đến UC-08, Giám sát viên UC-09 đến UC-11, Quản lý UC-12 đến UC-14, không use case nào dùng chung.
    \item \textbf{Quan hệ mở rộng.} UC-10 mở rộng UC-09 và UC-11 mở rộng UC-10 (tìm kiếm, xem một kết quả, lưu vào vụ việc); tương tự, UC-13 mở rộng UC-12 và UC-14 mở rộng UC-13. Các use case mở rộng vẫn nối với tác nhân, vì tác nhân có thể bắt đầu chúng độc lập, chẳng hạn mở danh sách vụ việc mà không tìm kiếm.
    \item \textbf{Đăng nhập là tiền điều kiện} của mọi use case khác, không phải một bước bên trong, nên không vẽ quan hệ «include» tới UC-01.
\end{itemize}

Các chức năng tự động (Mục~\ref{sec:4.2.5}) không có use case riêng vì không do tác nhân nào kích hoạt; chúng chạy khi xử lý AI được bật trong UC-04.

\begin{table}[H]
\centering
\caption{Danh sách use case của hệ thống}
\label{tab:usecases}
\begin{tabularx}{\textwidth}{|>{\raggedright\arraybackslash}p{1.4cm}|>{\raggedright\arraybackslash}p{3.8cm}|>{\raggedright\arraybackslash}p{2.4cm}|L|}
\hline
\textbf{Mã} & \textbf{Tên use case} & \textbf{Tác nhân} & \textbf{Mục đích} \\
\hline
UC-01 & Đăng nhập hệ thống & Cả ba tác nhân & Xác thực và truy cập các chức năng theo vai trò \\
\hline
UC-02 & Quản lý tài khoản người dùng & Quản trị viên & Tạo, cập nhật, khóa, xóa tài khoản; gán vai trò và khu vực \\
\hline
UC-03 & Quản lý camera & Quản trị viên & Thêm, cập nhật camera, kiểm tra kết nối RTSP, loại khỏi và đưa lại vận hành \\
\hline
UC-04 & Quản lý xử lý AI trên camera & Quản trị viên & Bật, tắt xử lý AI cho từng camera; xử lý tệp video dự phòng \\
\hline
UC-05 & Cấu hình mô hình AI & Quản trị viên & Chọn Detector và Tracker áp dụng cho toàn hệ thống \\
\hline
UC-06 & Theo dõi trạng thái hệ thống & Quản trị viên & Theo dõi trạng thái camera, luồng RTSP và xử lý AI \\
\hline
UC-07 & Kiểm tra hoạt động của AI & Quản trị viên & Kiểm tra nhóm xử lý camera và nhóm tìm kiếm \\
\hline
UC-08 & Xem nhật ký hệ thống & Quản trị viên & Tra cứu nhật ký thao tác và sự kiện kỹ thuật \\
\hline
UC-09 & Tìm kiếm người bằng AI & Giám sát viên & Tìm người bằng ảnh, câu mô tả hoặc thuộc tính trong khu vực được gán \\
\hline
UC-10 & Xem và đánh giá kết quả tìm kiếm & Giám sát viên & Xem chi tiết, đối chiếu và đánh giá từng kết quả \\
\hline
UC-11 & Quản lý hồ sơ vụ việc & Giám sát viên & Tạo vụ việc, lưu kết quả, cập nhật, đánh dấu hoàn thành, mở lại \\
\hline
UC-12 & Xem thông tin tổng quan & Quản lý & Xem các chỉ số và vụ việc gần đây trên toàn hệ thống \\
\hline
UC-13 & Xem hồ sơ vụ việc & Quản lý & Xem mọi vụ việc ở chế độ chỉ đọc \\
\hline
UC-14 & Xem thông tin kết quả tìm kiếm đã lưu & Quản lý & Xem chi tiết kết quả đã lưu trong vụ việc \\
\hline
UC-15 & Đăng xuất hệ thống & Cả ba tác nhân & Kết thúc phiên làm việc \\
\hline
\end{tabularx}
\end{table}
